\documentclass[aspectratio=169,10pt,spanish]{beamer}

\input{../../../_preambulo_beamer.tex}
\input{../../../_comandos_beamer.tex}

\selectlanguage{spanish}

\title{MA1001B - Modelación Estadística para la Toma de Decisiones}
\subtitle{Lección 31: Chi-cuadrada e intervalo de confianza para la varianza}
\author{}
\institute{Tecnol\'ogico de Monterrey}
\date{}

\begin{document}

\begin{frame}[plain]
  \vspace{1.2cm}
  {\Large\bfseries MA1001B - Modelación Estadística para la Toma de Decisiones\par}
  \vspace{0.7cm}
  {\large Lección 31: Chi-cuadrada e intervalo de confianza para la varianza\par}
  \vspace{0.7cm}
  {\color{TecRojo}\rule{\textwidth}{1.2pt}}
  \vspace{0.8cm}

  Facultad MA1001B\\[0.3cm]
  Periodo\\[0.3cm]
  Tecnol\'ogico de Monterrey
\end{frame}

\begin{frame}{La pregunta del intervalo para la varianza}
  \begin{alertblock}{Pregunta guía}
    ¿Cómo se construye un intervalo al 95\% para una $\sigma^{2}$ desconocida
    a partir de $s^{2}$ cuando la población se trata como normal?
  \end{alertblock}

  \vspace{0.6em}
  Al final de esta lección, usted puede:
  \begin{itemize}
    \item calcular $s^{2}$ con divisor $n-1$;
    \item obtener valores críticos chi-cuadrada con $\mathrm{gl}=n-1$;
    \item formar $((n-1)s^{2}/\chi^{2}_{U},(n-1)s^{2}/\chi^{2}_{L})$;
    \item explicar por qué los datos no normales invalidan el intervalo.
  \end{itemize}

  \vfill
  \scriptsize Todos los tiempos de ciclo usados hoy son sintéticos. No se usan datos reales de empresa.
\end{frame}

\begin{frame}{Una muestra normal pequeña de tiempos de ciclo}
  \small
  \begin{center}
    \begin{tabular}{lr}
      \toprule
      \textbf{Cantidad} & \textbf{Valor} \\
      \midrule
      Tamaño de muestra $n$ & 20 ciclos \\
      Grados de libertad $n-1$ & 19 \\
      $s^{2}$ muestral (semilla 42) & 6.814885 \\
      $\sigma^{2}$ oculta & 9 \\
      \bottomrule
    \end{tabular}
  \end{center}

  \vspace{0.5em}
  \begin{exampleblock}{Pivote}
    Si la población es normal, $(n-1)s^{2}/\sigma^{2}\sim\chi^{2}_{n-1}$.
    Ese pivote se invierte para obtener un intervalo para $\sigma^{2}$.
  \end{exampleblock}
\end{frame}

\begin{frame}[fragile]{Paso 1: Varianza muestral}
  \begin{columns}[T]
    \begin{column}{0.42\textwidth}
      \small
      \begin{lstlisting}
s2 = np.var(sample, ddof=1)
df = n - 1
      \end{lstlisting}

      \begin{block}{Salida verificada}
        $\mathrm{gl}=19$\\
        $s=2.610533$\\
        $s^{2}=6.814885$
      \end{block}
    \end{column}
    \begin{column}{0.58\textwidth}
      \centering
      \includegraphics[width=\textwidth]{../figuras/chi_cuadrada_01_varianza_muestral.png}
      \vspace{0.2em}
      \scriptsize Histograma de tiempos de ciclo del Paso 1
    \end{column}
  \end{columns}
\end{frame}

\begin{frame}{Invertir el pivote chi-cuadrada}
  \small
  Valores críticos de \texttt{scipy.stats.chi2.ppf}:
  \[
    \chi^{2}_{0.025,19}=8.906516,\qquad
    \chi^{2}_{0.975,19}=32.852327.
  \]
  Intervalo:
  \[
    \left[\frac{19\times 6.814885}{32.852327},
          \frac{19\times 6.814885}{8.906516}\right]
    =[3.941359, 14.537986].
  \]
  El valor chi-cuadrada más grande produce el límite inferior. Este intervalo
  con semilla 42 cubre la $\sigma^{2}=9$ oculta.
\end{frame}

\begin{frame}[fragile]{Paso 2: El intervalo chi-cuadrada al 95\%}
  \begin{columns}[T]
    \begin{column}{0.40\textwidth}
      \small
      \begin{lstlisting}
lower = df * s2 / chi2_u
upper = df * s2 / chi2_l
      \end{lstlisting}

      \begin{block}{Salida verificada}
        $\chi^{2}_{L}=8.906516$\\
        $\chi^{2}_{U}=32.852327$\\
        IC $[3.941359, 14.537986]$
      \end{block}
    \end{column}
    \begin{column}{0.60\textwidth}
      \centering
      \includegraphics[width=\textwidth]{../figuras/chi_cuadrada_02_intervalo_varianza.png}
      \vspace{0.2em}
      \scriptsize Intervalo para la varianza del Paso 2
    \end{column}
  \end{columns}
\end{frame}

\begin{frame}{Paso 3: Los datos no normales invalidan el intervalo}
  \begin{columns}[T]
    \begin{column}{0.42\textwidth}
      \small
      5000 replicaciones, $n=20$.

      \vspace{0.4em}
      Ciclos normales: cobertura $0.949400$.

      \vspace{0.4em}
      Ciclos exponenciales: cobertura $0.737200$.

      \vspace{0.5em}
      \textbf{Límite:} el intervalo chi-cuadrada para $\sigma^{2}$ supone una
      población normal. Tiempos de ciclo sesgados destruyen el 95\% nominal.
    \end{column}
    \begin{column}{0.58\textwidth}
      \centering
      \includegraphics[width=\textwidth]{../figuras/chi_cuadrada_03_limite_no_normal.png}
      \vspace{0.2em}
      \scriptsize Cobertura empírica del Paso 3
    \end{column}
  \end{columns}
\end{frame}

\begin{frame}{Verificación de conceptos}
  \begin{enumerate}
    \item ¿Por qué $s^{2}$ usa divisor $n-1$ y no $n$?
    \item ¿Qué cuantil chi-cuadrada va en el denominador del límite inferior?
    \item ¿El intervalo es simétrico alrededor de $s^{2}=6.814885$?
    \item ¿Qué ocurre con la cobertura cuando los tiempos de ciclo son
      exponenciales?
  \end{enumerate}

  \vspace{0.6em}
  \begin{block}{Conexión con la Lección 32}
    La sesión puede continuar directamente con la distribución $F$ y una
    comparación de dos varianzas.
  \end{block}

  \vfill
  \scriptsize Fuente: OpenStax, \textit{Introductory Business Statistics 2e},
  Capítulo 11, Sección 11.2. CC BY-NC-SA.
\end{frame}

\appendix



\end{document}
